\section{Background}
\label{sec:bg}

\subsection{The legal driver}
\label{sec:legal}

The \textbf{European Accessibility Act} (Directive (EU) 2019/882) requires products sold in the EU to
meet common accessibility requirements; every member state transposed it into national law by
June~2022. Its technical yardstick is \textbf{WCAG~2.2}. WCAG targets the web, so for a native app we
apply it through \emph{WCAG2ICT} and mobile guidance --- mapping each success criterion to its
Android equivalent rather than claiming literal conformance.

\subsection{The POUR principles}

WCAG rests on four principles (Figure~\ref{fig:pour}). For a media player they read concretely as:

\begin{itemize}[leftmargin=1.4em,itemsep=1pt]
	\item \textbf{Perceivable} --- label every transport icon; keep duration text at $\geq$4.5:1 contrast.
	\item \textbf{Operable} --- play, seek, and bookmark reachable from TalkBack, not only by touch.
	\item \textbf{Understandable} --- the play button announces ``playing'' vs.\ ``paused''.
	\item \textbf{Robust} --- expose a correct semantics tree instead of relying on pixels.
\end{itemize}

\begin{figure}[h]
	\centering
	\begin{tikzpicture}[
			box/.style={rounded corners=6pt,draw=brandblue,line width=1pt,fill=softbg,
					text width=3.5cm,minimum height=3.4cm,align=center,inner sep=6pt},
		]
		\def\cap#1{{\large\bfseries\color{brandblue}#1}}
		\node[box] (p) at (0,0) {\cap{P}erceivable\\[2pt]\small Equal information regardless of ability, via text alternatives and adjustable presentation.};
		\node[box] (o) at (4.2,0) {\cap{O}perable\\[2pt]\small Every action reachable by any input: touch, keyboard, switch, screen reader.};
		\node[box] (u) at (8.4,0) {\cap{U}nderstandable\\[2pt]\small Controls are predictable and clearly labelled, with helpful error feedback.};
		\node[box] (r) at (12.6,0) {\cap{R}obust\\[2pt]\small Reliable, consistent semantics that any assistive tech can interpret.};
	\end{tikzpicture}
	\caption{The four POUR principles that underpin WCAG.}
	\label{fig:pour}
\end{figure}

\subsection{Two trees, one screen}

The decisive fact for an Android developer: what is drawn and what a screen reader traverses are
\emph{two separate trees} (Figure~\ref{fig:trees}).

\begin{itemize}[leftmargin=1.4em,itemsep=1pt]
	\item \textbf{Composition tree} --- the visual hierarchy; nodes are boxes, text, and images ordered
	      by draw position.
	\item \textbf{Semantics tree} --- runs in parallel; each node carries a name, role, state, and
	      supported actions. It is the only thing assistive services can read.
\end{itemize}

\begin{figure}[h]
	\centering
	\begin{tikzpicture}[
		n/.style={rounded corners=4pt,draw,line width=0.8pt,font=\small,inner sep=5pt,align=center},
		vis/.style={n,draw=slate,fill=softbg},
		sem/.style={n,draw=brandgreen,fill=brandgreen!8},
		e/.style={-{Stealth[length=2mm]},draw=slate},
		se/.style={-{Stealth[length=2mm]},draw=brandgreen}
		]
		% Visual tree
		\node[vis] (root) at (0,0) {Card (Box)};
		\node[vis] (img)  at (-2.4,-1.5) {Image};
		\node[vis] (col)  at (0,-1.5) {Column};
		\node[vis] (btn)  at (2.4,-1.5) {IconButton};
		\node[vis] (t1)   at (-1.9,-3.1) {Text\\``Built Different''};
		\node[vis] (t2)   at (1.9,-3.1) {Text\\``Trung Kien Le''};
		\draw[e] (root)--(img); \draw[e] (root)--(col); \draw[e] (root)--(btn);
		\draw[e] (col)--(t1); \draw[e] (col)--(t2);
		\node[font=\bfseries\color{slate}] at (0,0.9) {Composition (visual) tree};

		% Semantics tree
		\begin{scope}[xshift=8.6cm]
			\node[sem] (sroot) at (0,0) {Merged node\\{\scriptsize name = ``Built Different, Trung Kien Le''}\\{\scriptsize role = Button, actions = \{Open, Bookmark\}}};
			\node[sem] (s1) at (-1.6,-2.2) {name:\\``Built Different''};
			\node[sem] (s2) at (1.6,-2.2) {name:\\``Trung Kien Le''};
			\draw[se] (sroot)--(s1); \draw[se] (sroot)--(s2);
			\node[font=\bfseries\color{brandgreen}] at (0,0.95) {Semantics tree};
		\end{scope}

		\draw[-{Stealth[length=2.5mm]},brandblue,line width=1pt,dashed]
		(3.0,-1.5) to[bend left=12] node[above,font=\scriptsize\color{brandblue}]{TalkBack reads} (5.7,-0.6);
	\end{tikzpicture}
	\caption{The visual tree (left) is drawn; the semantics tree (right) is read by TalkBack. The
		decorative image is dropped and the two texts are \emph{merged} into one node --- the techniques
		of §\ref{sec:core} and~§\ref{sec:adv}.}
	\label{fig:trees}
\end{figure}

TalkBack and Switch Access run as system services that subscribe to the \emph{merged} semantics tree,
turn each node into speech or a scan stop, and translate gestures back into actions. A pixel-perfect
screen is unusable if its semantics tree is missing, mislabelled, or misordered. Every technique in
this report shapes that tree.

\subsection{Jetpack Compose}

Compose is Android's declarative UI toolkit: the UI is a function of state, and state changes trigger
recomposition. It builds the semantics tree from Material components' built-in semantics plus
\code{Modifier.semantics} blocks --- giving code-level control over exactly what assistive technology
perceives.
